\documentclass[a4paper,11pt,fleqn]{article}%fleqn pour aligner les equations à gauche
\usepackage{../../Styles/Cours}

\newcommand{\chnum}{Chapitre 7}
\newcommand{\chtitle}{Force des acides}
\newcommand{\dtype}{TP}

  
\begin{document}
\maketitle

Des solutions d’acide éthanoïque comme le vinaigre ou d’acide chlorhydrique peuvent être utilisées pour détartrer des récipients en verre, de la robinetterie ou des appareils électroménagers. Les ions oxonium \ce{H3O+} libérés par les acides réagissent avec le tartre et le solubilisent. L’acide éthanoïque et l’acide chlorhydrique sont-ils aussi efficaces pour détartrer ? \\
\textbf{Problématique : Comment comparer la force de différents acides dans l’eau ?}\\
On comparera l’efficacité de l’acide ascorbique \ce{C6H8O6} (vitamine C), de l’acide chlorhydrique \ce{HCl} et de l’acide éthanoïque \ce{C2H4O2}. 

\noindent
\fbox{
\begin{minipage}{\textwidth}
\textbf{Document 1 : Matériel disponible}
\begin{multicols}{2}
\begin{itemize*}
\item acide ascorbique en poudre
\item solution d'acide chlorhydrique $S_1$ de concentration $c_1$ = \num{1,0e-2} \unit{\mole \per \liter}
\item solution d'acide éthanoïque $S_2$ de concentration $c_2$ = \num{1,0e-2} \unit{\mole \per \liter}
	\item pissette d’eau distillée
	\item balance
	\item spatule
	\item coupelle ou sabot de pesée
	\item pH-mètre et solutions étalon
	\item papier filtre
	\item béchers
	\item bécher poubelle
	\item éprouvette graduée de \num{50} \unit{\milli \liter}
	\item fiole jaugée de \num{50} \unit{\milli \liter}
\end{itemize*}
\end{multicols}
\end{minipage}}
\noindent
\fbox{\begin{minipage}{\textwidth}
\textbf{Document 2 : Données}\\
Masse molaire de l'acide ascorbique : $M_{AcA}$ = \num{176} \unit{\gram \per \mole}
\end{minipage}}

\section{Étude préalable}
\begin{enumerate}
	\item Proposer un moyen pratique pour connaître l'efficacité  d'un acide, par exemple pour jouer le rôle de détartrant
	\item Dans le tableau ci-dessous, écrire l'équation de réaction d'un acide \ce{AH} avec de l'eau.
	\item Compléter le tableau d'avancement de la réaction
\end{enumerate}

\begin{tabular}{|>{\centering}m{4cm}|>{\centering}m{2cm}|>{\centering}m{2cm}|>{\centering}m{2cm}||>{\centering}m{2cm}|>{\centering\arraybackslash}m{2cm}|}
\hline \multicolumn{2}{|c|}{} & \multicolumn{4}{c|}{\rule[1cm]{0cm}{0cm}}\\
\hline État initial & $x=0$ & & & & \rule[1cm]{0cm}{0cm}\\
\hline État final (si transformation totale) & $x_{max}$ & & & & \\
\hline
\end{tabular}
\begin{enumerate}[resume]
	\item En déduire l'expression de la concentration maximale attendue en ions oxonium notée $[\ce{H3O+}]_{max}$ en fonction de la quantité initiale d'acide $n_i(\ce{AH})$, puis en fonction de la concentration initiale d'acide $[\ce{AH}]_i$
\end{enumerate}

\section{Pratique}
	\subsection{Réaction d'un acide avec de l'eau}
	On fait réagir l'acide ascorbique avec de l'eau afin de préparer une solution $S_1$.
	\begin{enumerate}[resume]
		\item Préparer \num{50} \unit{\milli \liter} de solution $S_1$, en dissolvant \num{0,09} \unit{\gram} d'acide ascorbique. Déterminer la concentration initiale $c_1$ en acide ascorbique de la solution.
	\end{enumerate}
	Pour des questions de temps et de mise en œuvre, les réactions entre l'acide chlorhydrique et l'eau, d'une part , ont déjà été réalisées. Des solutions aqueuses de chaque acide sont donc à disposition. Pour garantir la validité de la comparaison, les trois solutions ont la même concentration initiale en acide.
	\subsection{Acide fort et acide faible}
	\begin{enumerate}[resume]
		\item À l’aide de la question 4, déterminer la concentration maximale théorique en ions oxonium $[\ce{H3O+}]_{max}$ de chaque solution puis inscrire les valeurs dans le tableau ci-dessous. 
	\end{enumerate}
	\begin{tabular}{|>{\centering}m{5cm}|>{\centering}m{4.2cm}|>{\centering}m{4.2cm}|>{\centering\arraybackslash}m{4.1cm}|}
	\hline & Solution $S_1$ d'acide ascorbique & Solution $S_2$ d'acide chlorhydrique & Solution $S_3$ d'acide éthanoïque\\
	\hline $[\ce{H3O+}]_{max}$ (en \unit{\mole \per \liter}) & & & \rule[1cm]{0cm}{0cm}\\
	\hline pH mesuré & & & \rule[1cm]{0cm}{0cm}\\
	\hline $[\ce{H3O+}]_{f}$ (en \unit{\mole \per \liter}) & & & \rule[1cm]{0cm}{0cm}\\
	\hline
	\end{tabular}
	\begin{enumerate}[resume]
		\item Mesurer le pH des trois solutions, puis déterminer leur concentration finale réelle $[\ce{H3O+}]_f$ en ions oxonium. 
		\item Pour chaque solution, comparer la concentration maximale attendue avec la concentration réelle. Que peut-on en déduire sur la réaction de chaque acide avec l’eau ?
		\item Écrire alors les équations de réaction de chaque acide avec l’eau de manière appropriée
		\end{enumerate}
L’un des trois acides est un \textbf{acide fort}, les deux autres sont des \textbf{acides faibles}.
	\begin{enumerate}[resume]
	\item Proposer une définition pour chacun de ces termes et expliquer lequel des trois acides serait le plus efficace pour détartrer. 
	\end{enumerate}
	\subsection{Constante d'acidité des acides faibles}
	Dans cette partie, on s'intéresse uniquement aux acides faibles. La force d'un acide faible peut être estimée à l'aide de la constante d'équilibre de la réaction entre l'acide et l'eau. Dans ce cas, la constante d'équilibre est alors nommée \textbf{constante d'acidité} $K_A$. Elle est caractéristique du couple acide-base étudié.
	\begin{enumerate}[resume]
	\item Donner l'expression de la constante d'acidité pour la réaction entre un acide faible et l'eau
	\item Déterminer sa valeur expérimentale pour chaque acide faible étudié
	\item À l'aide des questions de la partie précédente, identifier le lien entre la constante d'acidité $K_A$ et la force d'un acide faible
	\item Classer les trois acides par force croissante, répondre à la problématique
	\end{enumerate}
	
\end{document}